\documentclass[12pt,letterpaper,boxed]{hmcpset}
\usepackage[margin=1in,headheight=14pt]{geometry}
\usepackage{amsfonts, amsmath, amssymb, enumerate, fancyhdr, gensymb, lastpage, mathtools, parskip, graphicx}
\usepackage{xcolor, tikz-cd}
\newcommand{\wg}[1]{\textcolor{violet}{#1}}
\newcommand{\OO}{\mathcal O}
\newcommand{\Q}{\mathbb Q}
\newcommand{\R}{\mathbb R}
\newcommand{\C}{\mathbb C}
\newcommand{\Z}{\mathbb Z}
\newcommand{\abs}[1]{\left|#1\right|}
\newcommand{\im}{\text{im }}
\newcommand{\inv}{^{-1}}
\newcommand{\normal}{\unlhd} %% one can also use \trianglelelefteq
\newcommand{\anglee}[1]{\langle #1 \rangle}
\usepackage[shortlabels]{enumitem}

\newcommand{\pdv}[2]{\frac{\partial #1}{\partial #2}}

% Numbering macros
\pagestyle{fancy}
\lhead{Will Gilroy}
\chead{Complex Homework \#5}
\rhead{25 February 2026}
\lfoot{}
\cfoot{}
\rfoot{Page\ \thepage\ of\ \pageref{LastPage}}

\linespread{1.5}

\newcommand\blankpage{
    \thispagestyle{empty}
    \addtocounter{page}{-1}
    \newpage}
\renewcommand\footrulewidth{0.4pt}

\begin{document}

\problemlist{Complex Analysis Homework \#5} 

%------------------------- Problem 1 -----------------------

\begin{problem}
	\includegraphics[scale=0.8]{1.png}
	\hfill
\end{problem}

\begin{solution}
Consider the following. Let $U$ be the open set $\Delta \setminus \{0\}$
and let $u: U \to \R$ be the real logarithm $\ln(x^2 +y^2)$. Notice
that this function is indeed harmonic, we have that \[
	\frac{\partial^2 u}{\partial x^2} + \frac{\partial^2 u}{\partial x^2}
= \frac{2(y^2 - x^2)}{(x^2 + y^2)^2} + \frac{2(x^2 - y^2)}{x^2 + y^2}
= 0. 
\]
We show that this cannot be the real part of any holomorphic function
on $U$. Suppose $f: U \to \C$ is holomorphic with $f(x,y) = u(x,y) +
iv(x,y)$. We can write a condition on $f'$ in terms of $u$ using the
Cauchy-Riemann equations. Recall that we can write the complex
differential operator $\partial / \partial z = \frac{1}{2}(\pdv{}{x} +
\frac{1}{i}\pdv{}{y})$. And so we have
\begin{align*}
	\pdv{f}{z} 
	&= \frac{1}{2} \left( \pdv{}{x} + \frac{1}{i} \pdv{}{y} \right) (u + iv) \\
	&= \frac{1}{2} \left( u_x -i u_y + i v_x + v_y \right) \\
	&= \frac{1}{2} \left( u_x -i u_y + i (-u_y) + v_x \right) && \text{Cauchy-Riemann}\\
	f'(z) &= u_x - iu_y. 
\end{align*}
Applying this with $u(x,y) = \ln(x^2 + y^2)$ gives \[
	f'(z:=x+iy) = \frac{2(x- iy)}{(x^2 + y^2)^2} = \frac{2}{z}.
\]
Where the last equality comes from noticing \[
	\frac{1}{z} = \frac{1}{x + iy} = \frac{x-iy}{x^2 + y^2}. 
\]
I digress; we have shown that if $u = \ln(x^2 + y^2)$ were the
	real part of some holomorphic function $f$ then we would have
	$f'(z) = 2/z$. In other words, $f$ would need to be the
	antiderivate to $2/z$ on $\Delta \setminus \{0\}$. 
	However we have seen that $1/z$ has no antiderivative on any open
	set containing a circle centered at the origin. And so no such $f$
	exists on $\Delta \setminus \{0\}$. 
\end{solution}

\newpage

%------------------------- Problem 2 -----------------------

\begin{problem}
	\includegraphics[scale=0.8]{2.png}
	\hfill
\end{problem}

\begin{solution}

Let us denote our rational function $f(z) = P(z)/ Q(z)$ for complex
polynomials $P,Q$.
Suppose $f$ satisfies $\abs z = 1$ implies $\abs{f(z)} = 1$, then we
want to find conditions on $P,Q$.
Recall that $1 = \abs{z} = z \cdot \bar z$, i.e., $\bar z = 1/z$. Moreover,
$\abs{f} = 1$ when $\abs z = 1$ means that we have \[
	\frac{P(z)}{Q(z)} = \frac{\bar Q(1/z)}{\bar P(1/z)},
\]
where $\bar P, \bar Q$ mean the polynomial whose coefficients are the
complex conjugates of the coefficients of $P,Q$ respectively.
This shows that if $a$ is a root of $f$ then $1/a$ is a root of $\bar
f$. Moreover, if $b$ is a pole of $f$ then $1/b$ is a pole of $\bar
f$. Notice that $g:= \bar f$ is also a rational function (multiply $(\bar
Q(1/z))/(\bar P(1/z))$ by some appropriate power of $z^k$ with $k \in
\Z$) and so we have rational functions $f,g$ such that $fg = 1$
everywhere on $S^1$, which definitely contains a limit point as a
subset of $\C$. And so in fact $fg$ is the constant function $1$
everywhere on $\C$. Factoring $P,Q$ gives that we must have  \[
1 = f(z)g(z) = f(z) \bar f(1/z) = 
\frac{\Pi_k (z - a_k) \Pi_j (1/z -\bar b_j)}{ \Pi_j (z - b_j) \Pi_k(1/z - \bar a_k)}
= z^{\tilde m} \frac{\Pi_k (z - a_k) \Pi_j (1 - z\bar b_j)}{ \Pi_j (z
- b_j) \Pi_k(1 - z\bar a_k)},
\]
for some appropriate $\tilde m \in \Z$. 
\wg{I'll be honest, I'm not totally sure what to do with this.
Assuming that $P,Q$ have no common divisors then we need that each
factor $(z -a_k)$ must in fact equal some $(1 - z \bar a_k)$ and
likewise each factor $(z-b)$ must in fact equal some $(1 - z\bar b_j)$
(ish, I'm not so sure what to do with this $z^{\tilde m}$ out front.)}


\iffalse
Let us denote our rational function
$f(z) = (a_0 + a_1z + \cdots + a_n z^n)/ (b_0 + b_1z + \cdots b_m
z^m)$
We want to find conditions on $a_i, b_i$ such that $f(S^1) \subseteq
S^1$. By definition, $f$ preserves $S^1$ if $\abs{z} = 1$ implies
$\abs{f(z)} =1$.
Notice that we have \[
\abs{f(z)} = \frac{\sum_{i=0}^{n} \abs{a_i} \abs{z}^i}{\sum_{j=0}^{m}
\abs{b_j}\abs{z}^j}. 
\]
And so, if $\abs{z} = 1$ then we have \[
\abs{f(z)} = \frac{\abs{a_0} + \abs{a_1} \cdots + \abs{a_n}}{\abs{b_0} +
\abs{b_1} + \cdots + \abs{b_m}}
\]
\fi

\end{solution}

\newpage

%------------------------- Problem 3 -----------------------

\begin{problem}
	\includegraphics[scale=0.8]{3.png}
	\hfill
\end{problem}
\begin{solution}
Since $c \in S^1$ we have\footnote
{
Recall $c \in S^1$ means $c \cdot \bar c = \abs c = 1$. In other words
$\bar c = 1/c$. 
} 
$\bar c = 1/c$. And so if $p$ is a polynomial we want to show
that there exists some $c \in S^1$ such that $\abs{1/c - p(c)} \geq
1$. If $z \in S^1$ then $\abs{1/z - p(z)} = \abs{z} \abs{1/z - p(z)} =
\abs{1 - zp(z)}$. Let us denote $f(z) := 1 - zp(z)$. Notice that we
have $f(0) = 1$.

Suppose for contradiction that we strictly have $\abs{f(c)} < 1$ for all $c
\in S^1$. Since $f$ is non-constant and holomorphic on $\Delta$ we must have that $\abs{f(z)} < 1$
for all $z \in \Delta$ by the maximum modulus principle. 
Indeed, $\abs{f}: \overline \Delta \to \R$ is a continuous function
on a compact set and so by the extreme value theorem 
there elements in $\Delta$ such that $\abs{f}$
achieves its maximum and its minimum. If we had strictly that $\abs{f} < 1$
on $S^1$
but that $\abs{f(z_0)}
\geq 1$ for some $z_0 \in \Delta$ then that would mean $\abs{f}$ achieved its maximum for some
point not on $S^1$ (i.e. it's maximized on the interior $S^1$). 
Then the maximum modulus principle would conclude that $f$ is
constant. But $f(z) = 1 - zp(z)$ is nonconstant on $\Delta$ and so we
have that $\abs{f(z)} < 1$ for all $z \in \Delta$. 

We have that $\abs{f(z)} < 1$ on $S^1$ means that we must have
$\abs{f(z)} < 1$ for all $z \in S^1$. 
However, we
have that $f(0) = 1$, a contradiction. 
\end{solution}

\newpage

%------------------------- Problem 4 -----------------------

\begin{problem}
	\includegraphics[scale=0.8]{4.png}
	\hfill
\end{problem}

\begin{solution}
I must admit, I'm not suuuper sure how to do this problem, but here
are some thoughts. The only approximation theorem we've learned in
class is the Runge approximation theorem. However, that does not seem
to apply here since the desired $f$ is not continuous on any
neighbourhood of $0$, let alone holomorphic.

On the other hand a similar function \[
	g(z) := \begin{cases}
		-1 & Re(z) < 0  \\
		0 & Re(z) \geq 0
	\end{cases}
\]
is equal to $H(Re(z)) -1$ where $H$ is the heaviside function, or the
step function with a single step. We can approximate $H(x)$ pointwise as \[
	H(x) = \lim_{k \to \infty} \frac{1}{1 + e^{-2kx}}.
\]
Consider, for a fixed $k$, we then have \[
	g_k(z) 
	= \frac{1}{1 + (1 + (-2kRe(z) + (-2kRe(z))^2/2 \cdots + )} - 1 
= \frac{1 + (-2kRe(z)) + (-2kRe(z))^2/2  + \cdots}{2 + (-2kRe(z) +
(-2kRe(z))^2/2 + \cdots}.
\]
Recalling that $Re(z) = (z + \bar z)/2$ then the above expression is
at least a rational power series in $z, \bar z$ (although perhaps not
a rational function of $z$). We can truncate the power series in the
numerator and denominator to then have a rational power series
approximation of $g_k(z)$, which in turn is an approximation of $g$. 

Now, in terms of heaviside functions we can write the desired function
as \[
f(z) = (H(Re(z)) - 1) + H(Re(z))e^z \approx \frac{e^z -
e^{-2kRe(z)}}{1 + e^{-2kRe(z)}}. 
\]
Taylor expanding the numerator and denominator of this expression
would then give us some rational power series in $z, \bar z$ which we
could then truncate to get some rational polynomial (in $z, \bar z$)
approximation of the function we want. (However, this is rational and
not polynomial, and moreover is rational in $z, \bar z$ and not just
in $z$).
\end{solution}

\newpage

\end{document}
